\documentclass[12pt,a4paper]{article}
\usepackage[T1]{fontenc}
\usepackage[utf8]{inputenc}
\usepackage[english]{babel}
\usepackage{amsmath,amssymb,amsthm}
\usepackage{geometry}
\geometry{margin=2.5cm}
\usepackage{graphicx}
\usepackage{hyperref}
\hypersetup{colorlinks=true,linkcolor=blue,citecolor=blue,urlcolor=blue}
\usepackage{booktabs}
\usepackage{enumitem}

\newtheorem{theorem}{Theorem}
\newtheorem{remark}{Remark}

\title{\textbf{A Gauge Theory of Line Defects in an Elastic Continuum:\\
Dislocation Rings, the Peach--Koehler Force, and Emergent Interaction Laws\\
{\large (A Geometrodynamic Model of Particles as Elastic Defects)}}}
\author{A.~Yu.~Drozdov}
\date{24 October 2026}

\begin{document}
\maketitle

\begin{abstract}
\noindent
We present a consistent account of a geometrodynamic model in which elementary
particles are treated as topological defects (dislocation and disclination
rings) of an elastic aether, and the fundamental interactions appear as
different projections of a single elastic response of the medium. The model is
built on the Kadi\'c--Edelen and Lazar gauge theory of dislocations
\cite{lazar2010,kadic1983}, in which the magnetic field is identified with the
dislocation density and the Peach--Koehler force plays the role of the Lorentz
force. Gravity is introduced through a ``Fermi film'' mechanism --- a
deformable boundary of the action integral held orthogonal to the curved world
tubes of accelerating charges (a generalization of Fermi's 1923 treatment of
the $4/3$ paradox \cite{fermi1923}). We show that the inverse-square law for
gravity is a \emph{theorem} (a consequence of the Laplacian Green function),
not a fit. The model is formulated in an elastic unit basis $(L,T,M)$, without
invoking electric charge as a primary entity. The limits are drawn honestly:
the \emph{structure} of the interactions (power laws, signs, hierarchy) is
derived, together with a set of relations between scales, whereas the absolute
values of the fundamental constants ($\alpha$, $G$) remain input parameters.
The paper is a working preprint intended as a first introduction to the
subject.

\medskip
\noindent\textbf{Keywords:} gauge theory of dislocations, aether, Fermi film,
magnetic monopole, Burgers vector, Peach--Koehler force, $4/3$ paradox,
geometrodynamics.
\end{abstract}

\tableofcontents
\newpage

\section{Introduction}

\subsection{Motivation: the particle as a defect of the medium}
The idea that elementary particles may be not point objects but extended
structures in some medium has a long history --- from Kelvin's vortices to
modern soliton models. The present work develops a concrete version of this
programme: \emph{a particle is a topological defect of an elastic continuum}
(the ``aether''), namely a closed dislocation line (ring) or its
generalizations. The approach is attractive because it rests on the
well-developed mathematical apparatus of the theory of crystal defects
\cite{hirth1982,kadic1983} and its gauge formulation \cite{lazar2010}.

\subsection{Why dislocations}
In crystals, dislocations are line defects characterized by the Burgers vector
$\mathbf{b}$, which measures the closure failure of a circuit around the defect
line. M.~Lazar \cite{lazar2010} showed that the dynamic gauge theory of the
translation group $T(3)$ leads to equations structurally identical to Maxwell's
equations (Table~\ref{tab:corr}). This provides a physically meaningful
``dictionary'' between defect theory and electrodynamics.

\begin{table}[h]
\centering
\caption{Correspondence Maxwell~$\leftrightarrow$~dislocation theory
(after Lazar \cite{lazar2010}).}
\label{tab:corr}
\begin{tabular}{ll}
\toprule
Electrodynamics & Dislocation theory \\
\midrule
Magnetic field $\mathbf{B}$ & dislocation density $T^a$ \\
Electric field $\mathbf{E}$ & dislocation current $I^a$ \\
Magnetic excitation $\mathbf{H}$ & pseudomoment stress $H_a$ \\
Magnetic flux $\oint\mathbf{B}\,d\mathbf{S}$ & Burgers vector $\oint T^a = b$ \\
Lorentz force & Peach--Koehler force \\
\bottomrule
\end{tabular}
\end{table}

\subsection{The key feature: the charge is magnetic, not electric}
In the proposed model \emph{electric charge is not a primary entity}. What is
primary is the magnetic charge, identified with the flux of the dislocation
density, i.e.\ with the Burgers vector. The electron is built as a precessing
magnetic dipole. This is consistent with the long-standing line connecting
dislocations and Dirac monopoles \cite{mistura1995}.

\subsection{Structure of the paper}
Section~\ref{sec:units} introduces the elastic unit basis.
Section~\ref{sec:particle} describes the particle as a precessing dipole.
Section~\ref{sec:fermi} presents the Fermi-film mechanism and the derivation of
gravity. Section~\ref{sec:channels} collects the six interaction channels.
Section~\ref{sec:limits} honestly delimits the model.

\section{The elastic unit basis}
\label{sec:units}
We build a system of units from three primitives of the elastic medium: length
$L$, time $T$ and mass $M$. There is \emph{no} electric charge and no ampere in
the basis. The main quantities are:
\begin{itemize}[nosep]
\item Burgers vector $b$ (lattice displacement) --- dimension $[L]$;
\item shear modulus $\mu$ --- dimension $[M L^{-1}T^{-2}]$ (Pa);
\item density $\rho$ --- $[M L^{-3}]$;
\item speed of sound/light $c=\sqrt{\mu/\rho}$ --- $[L T^{-1}]$.
\end{itemize}
The magnetic charge as the flux of the dislocation density
$q_m=\oint T^a\,dS$ has the dimension of length $[L]$ --- this is internally
consistent in the elastic theory (unlike in SI, where the magnetic charge is
measured in A$\cdot$m; transferring SI units to the elastic medium is
unnecessary and leads to confusion).

\section{The particle as a precessing dipole}
\label{sec:particle}

\subsection{Geometry}
The electron is modelled as a pair of dislocation rings (of radius $a$) whose
planes are \emph{fixed}, while along each ring a ``special point'' --- a
localized magnetic charge $q_m=\Delta b$ (a jump of the Burgers vector) ---
runs with frequency $\omega$. By ``precession'' we mean precisely the motion of
the special point along the fixed ring, not a wobbling of the planes. The two
rings are separated by a distance $d$ and carry opposite charges $\pm q_m$
(the $N$ and $S$ poles), forming a magnetic dipole.

\subsection{Mass from self-interaction}
Mass is not introduced separately but is the elastic self-interaction energy of
the defect:
\begin{equation}
m = \frac{\mu\, b^2 a}{c^2},
\end{equation}
which has the correct dimension $[M]$.

\subsection{Magnetic moment and spin}
The magnetic moment is a \emph{static} dipole moment $\mu_{\text{mag}}=q_m\,d$
(not a current-loop moment: the ring current is magnetic, and the naive
current-loop formula does not apply here). Spin is tied to topology: a pair of
ordinary rings has configuration period $2\pi$ (a boson), whereas a pair of
M\"obius rings (ring + anti-ring) has period $4\pi$, corresponding to a fermion
of spin $1/2$. In the M\"obius pair the two rings carry \emph{opposite}
directions of the Burgers vector ($+\Delta b$ and $-\Delta b$, the $N$ and $S$
poles) --- hence a static magnetic dipole. Crucially, the Burgers twist
$b(\varphi)=\cos(\varphi/2)$ is a Frank rotation by $\pi$ per circuit, i.e.\ a
\emph{disclination} ($SO(3)$) degree of freedom: the M\"obius sector of the
model already lies in $SO(3)$, and the magnetic moment $\mu=\nu=n/2$ follows as
the flux of the spin connection (Osipov \cite{osipov2010}). Hydrodynamically
the same object is a vortex: the vorticity corresponds to the Frank density,
and the precession ($\omega>0$) required for stability is the circulation of
that vortex.

\subsection{Stability and annihilation}
Opposite charges $\pm q_m$ attract and, without motion, would collapse.
Precession provides a centrifugal barrier (conservation of angular momentum
$J=m a^2\omega$) that stabilizes the pair at a finite distance:
\begin{equation}
E(d) = -\frac{A}{d} + \frac{J^2}{2 m d^2},\qquad d_{\text{eq}}=\frac{J^2}{mA}.
\end{equation}
As $\omega\to0$ the barrier vanishes --- the pair annihilates (Frenkel
recombination, the total Burgers vector going to zero).

\section{The Fermi film and the derivation of gravity}
\label{sec:fermi}

\subsection{Historical context: the $4/3$ paradox}
E.~Fermi in 1923 \cite{fermi1923} showed that the well-known discrepancy
between the electrodynamic ($\tfrac{4}{3}U/c^2$) and relativistic ($U/c^2$)
masses of a charged body is removed if the boundary of the action integral is
taken \emph{strictly orthogonal} to the curved world tube of the accelerating
charge. We call this deformable boundary the \emph{Fermi film}. The film is not a
term of the Lagrangian but the boundary of the integration domain of the
action. It is therefore not a dynamical field, is not retarded, and adjusts
instantaneously to the curvature of the world tubes. Only the field itself,
obtained by minimizing the action, is retarded and radiates. The Fermi film should be distinguished from the Gibbons--Hawking--York
boundary term of general relativity \cite{hehl1995}. Our film differs from it in
that it does not enter the Lagrangian: the Gibbons--Hawking--York term is a term
of the action and obeys the field dynamics, whereas the Fermi film is merely the
boundary of the integration domain. Therefore it is not retarded and adjusts
instantaneously to the curvature of the world tubes.

\subsection{The source of deformation}
An accelerating charge curves its world tube; the film, remaining orthogonal,
is deformed. The density of the deformation source is proportional to the
\emph{square} of the proper acceleration $W=\omega^2 a$:
\begin{equation}
\sigma = \mu\, b^2\,\frac{\langle W^2\rangle}{c^2} > 0 .
\end{equation}
The sign is always positive (a square), so gravity is always attractive and
does not depend on the sign of the charge (universality).

\subsection{The inverse-square law as a theorem}
\begin{theorem}[$1/R^2$ law]
If the film deformation $\Delta$ obeys the equilibrium equation of an elastic
membrane with a point source,
\begin{equation}
\gamma\,\nabla^2\Delta(\mathbf r) = -\sigma\,\delta^3(\mathbf r-\mathbf r_0),
\end{equation}
then the solution is the Laplacian Green function, $\Delta(R)=\sigma/(4\pi\gamma R)$,
whence the interaction potential $U(R)\propto 1/R$ and the force $F(R)\propto 1/R^2$.
\end{theorem}
The exponent follows from the three-dimensionality of space, not from a fit.

\subsection{Stiffness parameters: two scales}
\paragraph{What $\mu_{\text{aeth}}$ is.}
The quantity $\mu_{\text{aeth}}$ is the \emph{shear modulus of the vacuum}, the
stiffness of the aether with respect to shear deformation, of pressure
dimension $[\text{Pa}]$. It is through this elastic modulus that, in the
$(L,T,M)$ basis (without the ampere), the main quantities of the model are
defined:
\begin{align}
c &= \sqrt{\mu_{\text{aeth}}/\rho} && (\text{transverse-wave speed}),\\
m &= \mu_{\text{aeth}}\,b^2 a / c^2 && (\text{particle mass}),\\
F &= \mu_{\text{aeth}}\,b^2 / R \cdot (\text{geometry}) && (\text{Coulomb force}).
\end{align}
Physically $\mu_{\text{aeth}}$ governs shear deformations (dislocations), from
which charge and mass are built; hence it sets the physics at the particle
scale (of order the Compton length $\lambda_C$). Its numerical value is
obtained by \emph{calibration}: substituting the electron mass $m_e$ and its
classical radius $r_e$ into the self-energy of the dislocation ring and solving
for $\mu_{\text{aeth}}$. This is \emph{not} a derivation from first principles
but the fixing of a single number by the electron, which we state explicitly.

A check against the model shows that \emph{one} modulus cannot account for both
mass/Coulomb and gravity: the required values differ by $\sim\!84$ orders of
magnitude. Hence two scales are needed:
\begin{align}
\mu_{\text{aeth}} &\approx 2.7\times10^{28}~\text{Pa} && (\text{EM/mass, Compton scale}),\\
\gamma_P &= \frac{c^7}{8\pi\hbar G^2}\approx 1.8\times10^{112}~\text{Pa} && (\text{gravity, Planck scale}).
\end{align}
The gravitational constant $G$ enters $\gamma_P$: it is an \emph{input}
parameter of the aether stiffness, not derived from more fundamental quantities.

\paragraph{What $\gamma_P$ is, and the nature of the hierarchy.}
The Planck modulus is the Planck \emph{energy density} divided by $8\pi$:
\begin{equation}
\gamma_P=\frac{c^7}{8\pi\hbar G^2}=\frac{1}{8\pi}\,\frac{m_P c^2}{\ell_P^3},
\qquad \ell_P=\sqrt{\hbar G/c^3},\quad m_P=\sqrt{\hbar c/G}
\end{equation}
(numerical agreement $<1\%$). This is not an arbitrary combination but the
stiffness of the vacuum at the Planck scale; the factor $8\pi$ is the same as
in the Einstein equations. Likewise $\mu_{\text{aeth}}=m_e c^2/(r_e^2\lambda_C)$
is the energy density at the electron scale. The ratio of the two moduli
\begin{equation}
\frac{\gamma_P}{\mu_{\text{aeth}}}\sim\frac{m_P}{m_e}\ \text{(with a volume factor)}\sim10^{84}
\end{equation}
reduces to the ratio of the Planck and Compton scales, i.e.\ to the smallness
of $m_e/m_P$. Thus the hierarchy of the two stiffnesses is \emph{equivalent} to
the standard \emph{hierarchy problem} ($m_e\ll m_P$): the model does not solve
it, but neither does it introduce a new arbitrary parameter --- it inherits the
well-known smallness of $m_e/m_P$.

\section{The six interaction channels}
\label{sec:channels}
All interactions are projections of a single elastic response. Their decay laws
and sources are collected in Table~\ref{tab:channels}; the power laws are
structural consequences (Green function, dipole--dipole interaction, mediator
mass, confinement), not the result of a fit.

\begin{table}[h]
\centering
\caption{The six interaction channels of the model.}
\label{tab:channels}
\begin{tabular}{llll}
\toprule
Channel & Law & Source & Coupling \\
\midrule
Coulomb & $1/R^2$ & product of Burgers charges & $\alpha=7.30\times10^{-3}$ \\
Magnetic & $1/R^4$ & dipole--dipole $q_m d$ & $\sim\alpha/4$ \\
Gravity & $1/R^2$ & $\langle W^2\rangle$ of the Fermi film & $\alpha_G=1.75\times10^{-45}$ \\
Exchange & $1/R^3$ & precession phase $\cos2\Delta\phi$ & $\sim\alpha$ \\
Weak & $e^{-R/R_W}$ & twist chirality & $\sim 1/30$ \\
Strong & $\propto R$ & fractional $b/3$ + ribbon & $\sim 1$ \\
\bottomrule
\end{tabular}
\end{table}

The whole system closes on two dimensionless numbers:
\begin{equation}
\alpha = \frac{k e^2}{\hbar c}=\frac{1}{137.04},\qquad
\alpha_G = \left(\frac{m_e}{m_{\text{Planck}}}\right)^2 = 1.75\times10^{-45},
\end{equation}
whence the weakness of gravity relative to the Coulomb force:
$\alpha_G/\alpha\approx 2.4\times10^{-43}$.

\section{Limits of the model: what is proved, assumed, open}
\label{sec:limits}
In the spirit of scientific honesty we explicitly separate the levels of
certainty.

\paragraph{Proved (structurally).} The decay laws of all channels; the
attractiveness and universality of gravity; the dimensional closure of mass;
the consistency of the magnetic charge as a length-dimension quantity in the
elastic basis; the annihilation mechanism and the protection by precession; the
Frank vector $\Omega=\pi$ of the $\cos(\varphi/2)$ twist (a wedge
half-disclination, $\nu=1/2$, spin $1/2$).

\paragraph{Assumed (input parameters).} The constants $\alpha$ and $\alpha_G$;
the gravitational constant $G$ (embedded in the stiffness $\gamma_P$); the
existence of two stiffness scales and their ratio.

\paragraph{Open.} The origin of the numerical values of $\alpha$ and
$\alpha_G$; a rigorous derivation of the gyromagnetic ratio $g=2$ (in the
classical ring picture one gets $g=1$; $g=2$ requires a spinor level of
description in the spirit of \cite{osipov2010}); the dimensional factor in the
gravity source, which requires a dynamical (time-dependent) film stiffness; the
numerical values of the weak- and strong-sector couplings. The tensor ($J=2$)
channel of the trace-free quadrupole $Q_{ij}$ exists structurally (spin-2) but
gives only a short-range $1/R^5$ correction; the long-range $1/R^2$ gravity is
intrinsically scalar (the trace $\langle W^2\rangle$).

\section{Conclusion}
We have proposed an internally consistent model in which particles are elastic
defects of the aether, and the interactions are projections of a single medium
response through the Peach--Koehler force and the deformation of the Fermi
film. The model explains the \emph{structure} (laws, signs, hierarchy of
interactions) without claiming to derive the fundamental constants from first
principles. Such a formulation makes the theory falsifiable: for
magneto-gravitational devices, for example, it predicts a negligible
($\sim\!10^{-13}$) relative correction to weight, incompatible with claims of
macroscopic antigravity. Further work includes a rigorous spinor derivation of
$g=2$, a refinement of the dynamical film stiffness, and this English version
for archival publication.

\begin{thebibliography}{99}
\bibitem{fermi1923} E.~Fermi, \emph{Correzione di una contraddizione tra la teoria
elettrodinamica e quella relativistica delle masse elettromagnetiche},
Nuovo Cimento \textbf{25}, 159 (1923).
\bibitem{lazar2010} M.~Lazar, \emph{On the fundamentals of the three-dimensional
translation gauge theory of dislocations}, Math.\ Mech.\ Solids \textbf{16}(3), 253 (2010).
\bibitem{kadic1983} A.~Kadi\'c, D.~G.~B.~Edelen, \emph{A Gauge Theory of Dislocations
and Disclinations}, Lecture Notes in Physics \textbf{174}, Springer (1983).
\bibitem{lazar2009} M.~Lazar, \emph{On the Higgs mechanism and stress functions in the
translational gauge theory of dislocations}, arXiv:0903.0990 (2009).
\bibitem{lazar2009b} M.~Lazar, C.~Anastassiadis, \emph{The gauge theory of dislocations:
static solutions of screw and edge dislocations}, arXiv:0802.0670 (2009).
\bibitem{osipov2010} E.~A.~Kochetov, V.~A.~Osipov, \emph{Dirac fermions on a disclinated
flexible surface}, arXiv:1001.4450 (2010).
\bibitem{osipov1999} E.~A.~Kochetov, V.~A.~Osipov, \emph{Gauge theory of disclinations on
fluctuating elastic surfaces}, arXiv:math-ph/9903018 (1999).
\bibitem{hehl1995} F.~W.~Hehl, J.~D.~McCrea, E.~W.~Mielke, Y.~Ne'eman, \emph{Metric-affine
gauge theory of gravity}, Phys.\ Rep.\ \textbf{258}, 1 (1995); arXiv:gr-qc/9402012.
\bibitem{hirth1982} J.~P.~Hirth, J.~Lothe, \emph{Theory of Dislocations}, 2nd ed.,
Wiley (1982).
\bibitem{mistura1995} L.~Mistura, \emph{Crystal dislocations and Dirac monopoles},
Int.\ J.\ Eng.\ Sci.\ \textbf{33}, 2149 (1995).
\bibitem{rohrlich1960} F.~Rohrlich, \emph{Self-energy and stability of the classical
electron}, Am.\ J.\ Phys.\ \textbf{28}, 639 (1960).
\end{thebibliography}

\end{document}
